\chapter{Efficiency and the Cost of Six Metrics at Scale (Day 8)}
\label{chap:day08}

\section{Motivation}

Efficiency is one of the six metrics defined for evaluating black-box explainable-AI frameworks
\citep{arreche2024exai}: not whether an explanation is faithful or stable, but whether producing it
is computationally affordable at the scale a real deployment or a larger study would require.
Chapters 3 and 5--7 each ran Florence-2 inference for a different purpose --- baseline judgments,
masking-based reruns, cross-attention attribution, stochastic-decoding reruns --- but none of them
asked, in isolation, what the combined cost of running all four looks like, or whether that cost
would still be reasonable at ten times this pilot's 163-sample size. Day 8's task was to answer
exactly that question: aggregate the per-sample timing already implicit in the pilot's four
inference-bearing days, extrapolate to a deployment-relevant scale, and report honestly on any gaps
in the data needed to do so.

\section{An Instrumentation Gap, Handled Honestly}
\label{sec:day08-gap}

The plan called for an \texttt{aggregate\_timings()} routine to pull \texttt{inference\_ms} straight
out of Days 3, 5, 6, and 7's CSVs, with zero new model calls required. Checking the actual column
headers first, rather than assuming the plan's premise held, showed that premise was only half
true: of the four CSVs, only \texttt{baseline\_predictions.csv} (Day 3) ever logged an
\texttt{inference\_ms} column at all. \texttt{descriptive\_accuracy.csv} (Day 5),
\texttt{visual\_sparsity.csv} (Day 6), and \texttt{stability.csv} (Day 7) have no timing column
whatsoever. Pulling pre-logged numbers from three of the four days was simply not available as an
option, and the gap was handled differently depending on what each day's CSV could and could not
support, rather than papered over with a single uniform assumption.

\textbf{Day 5 required no new model calls.} Day 5's descriptive-accuracy reruns use the
\emph{exact same call type} Day 3 already timed: \texttt{answer\_rule()} at the default
\texttt{num\_beams=3}. Because the call type is identical, Day 5's cost can be recovered exactly
from data already on disk, with no new inference run --- as \texttt{rerun\_count} $\times$ that
sample's own Day 3 \texttt{inference\_ms}, where \texttt{rerun\_count} (1, or 2 if a second region
existed for that sample) comes directly out of \texttt{region\_extraction.csv}'s \texttt{n\_regions}
column. This is exact recovery, not estimation: every input to the computation is a real, already-
measured quantity from this pilot's own data.

\textbf{Days 6 and 7 required a small, bounded calibration.} Day 6's cross-attention attribution
and Day 7's stability reruns both use decoding configurations that nothing in the pilot had timed
before --- \texttt{output\_attentions=True} (Day 6, forcing Florence-2's slower eager-attention
path, as documented in Chapter~\ref{chap:day06}) and \texttt{do\_sample=True, num\_beams=1} (Day 7,
genuine stochastic sampling rather than beam search). Recovering their cost exactly from existing
logs was not possible, because no existing log contains a timing for either configuration. Rather
than invent a number or silently reuse Day 3's unrelated timing, this script ran one small, fresh
\textbf{15-sample calibration} for each configuration and scaled the result. This is a deliberate,
bounded deviation from the plan's original "zero new model calls" premise, scoped the same
disciplined way Chapter~\ref{chap:day06}'s own architecture deviation was scoped: small enough not
to threaten Day 8's timebox, and documented here explicitly rather than folded silently into the
results table as if it were exact.

\section{Results}

\texttt{scripts/08\_efficiency\_analysis.py} assembled Table~\ref{tab:day08-results} from the
sources described in Section~\ref{sec:day08-gap}: exact aggregation for Days 3 and 5, calibrated
extrapolation for Days 6 and 7. Day 4 (region extraction) and the cross-day analysis carried out in
Chapter~11 perform no model inference at all --- both are pandas/PIL operations over CSVs Days 3 and
6 already produced --- and are correctly excluded from the table rather than padded in at zero cost.

\begin{center}
\begin{tabular}{p{1.1cm}p{3.6cm}p{4.1cm}rr}
\toprule
\textbf{Day} & \textbf{What} & \textbf{Calls/sample} & \textbf{Mean ms/sample} & \textbf{Extrap.\ to $n=$1,000} \\
\midrule
3 & Baseline inference & 1 \texttt{answer\_rule} call (2 generate calls inside) & 271 & 4.5 min \\
5 & Descriptive-accuracy reruns & 1.94 \texttt{answer\_rule} calls & 526 & 8.8 min \\
6 & Cross-attention attribution (15-sample calibration) & 1 \texttt{cross\-\_attention\-\_heatmap} call (1 generate call) & 215 & 3.6 min \\
7 & Stability sampled reruns & 3 \texttt{answer\_rule} calls & 778 & 13.0 min \\
\bottomrule
\end{tabular}
\captionof{table}{Per-sample inference cost by day, on the 163-sample pilot, linearly extrapolated to
$n=1{,}000$.}
\label{tab:day08-results}
\end{center}

Summed across all four rows, the \textbf{full six-metric pipeline at $n=1{,}000$, by linear
extrapolation, costs approximately 0.50 GPU-hours (about 30 minutes) on a single RTX 3070}.
Practically trivial: efficiency is conclusively not a blocker for scaling this pilot to the full
3,004-image ConstructionSite test split, let alone to a 1,000-sample study.

\subsection{Finding 1: comparing ``ms per call'' across days needs call-count normalizing}

Read at face value, Table~\ref{tab:day08-results} appears to contradict
Chapter~\ref{chap:day06}'s own finding that \texttt{output\_attentions=True} forces a slower
eager-attention path: Day 6's 215ms/sample looks \emph{cheaper} than Day 3's 271ms/sample, not more
expensive. The catch is a call-count mismatch hiding inside the per-sample average.
\texttt{answer\_rule()} (Days 3, 5, and 7) makes \textbf{two} \texttt{generate()} calls per logged
"call" --- one for the worker phrase, one for the safety-object phrase --- while
\texttt{cross\_attention\_heatmap()} (Day 6) makes only \textbf{one}. The two numbers are not
measuring the same unit of work, and comparing them directly is comparing apples to (half-)oranges.

Normalizing to a single-\texttt{generate()}-call basis removes the mismatch. Day 3's per-call cost
is its 271ms divided across its two internal generate calls, $271/2 \approx 135$ms; Day 7's
778ms/sample, divided across its three rerun calls (each itself two generate calls, so six generate
calls per sample), gives the same per-rerun-call figure as Day 3's, $\approx 259$ms per rerun call,
implying $259/2 \approx 130$ms per generate call --- consistent with Day 3's 135ms, as it should be,
since both use the same \texttt{num\_beams=3}-vs-sampling decoding intensity per generate call.
Day 6's single call, by contrast, costs the full 215ms on its own, with no division needed because
there is only one generate call inside it. Normalized this way, Day 6's per-generate-call cost is
\textbf{about 65\% more expensive} than Day 3 or Day 7's per-generate-call cost ($215$ vs.\
$\approx 130$--$135$ms) --- consistent with the eager-attention overhead Chapter~\ref{chap:day06}
already documented, once the comparison is made apples-to-apples. The lesson generalizes beyond this
one pair of rows: any cross-day efficiency comparison in this pilot has to account for how many
generate calls are bundled inside each logged "call" before treating raw per-sample milliseconds as
comparable.

\subsection{Finding 2: \texttt{num\_beams=3} beam search is barely slower than \texttt{num\_beams=1} sampling here}

A second, independent finding falls out of the same table. Day 7's sampled calls
(\texttt{num\_beams=1}, $\approx 259$ms per rerun call) are only about \textbf{4\% faster} than Day
3's beam-search calls (\texttt{num\_beams=3}, 271ms per call) --- not the roughly 3$\times$ gap a
naive "3 beams means 3$\times$ the decode compute" intuition would predict. The likely reason is
that Florence-2's grounding outputs in this pilot are very short: location-token sequences of only
about 7--8 generated tokens, per the token-count investigation already carried out in
Chapter~\ref{chap:day06}. With so few decode steps, each call's latency is dominated by the
\emph{fixed} cost of encoding the image through the DaViT vision tower once per call --- a cost that
does not depend on beam count at all --- rather than by the decode loop, where beam multiplicity
would actually matter. Beam count would likely matter far more for tasks that generate long
sequences, such as \texttt{<MORE\_DETAILED\_CAPTION>}; it barely registers for this pilot's short
grounding-token outputs.

\section{Worked Example}

Figure~\ref{fig:day08-worked} breaks the headline per-sample numbers in Table~\ref{tab:day08-results}
down visually, showing the relative cost of each day's inference workload behind the aggregate
figures reported above.

\begin{figure}[h]
\centering
\includegraphics[width=0.8\linewidth]{efficiency/per_sample_cost.png}
\caption{Per-sample inference cost by day (Days 3, 5, 6, 7), the timing breakdown behind
Table~\ref{tab:day08-results}'s headline numbers and the $n=1{,}000$ extrapolation.}
\label{fig:day08-worked}
\end{figure}

\section{Sanity Check on the Extrapolation}

Because Days 6 and 7's per-sample costs rest on a 15-sample calibration rather than the full
163-sample run, the extrapolation those two rows feed into was checked against an independent
order-of-magnitude estimate before being trusted. Day 3 and Day 5's numbers carry no such risk in
the first place: they are exact sums over the real 163 samples (44.2s and 85.7s respectively), with
no extrapolation step and therefore no extrapolation error possible. Day 6 and 7's numbers are a
15-sample calibration scaled linearly, first to the 163-sample pilot and then on to $n=1{,}000$, and
that scaling step is exactly where an error in the calibration would show up if one were present.

The calibrated 778ms/sample figure for Day 7 implies approximately \textbf{127 seconds} of pure GPU
compute for all $163 \times 3$ sampled reruns. That figure is the right order of magnitude against
what was actually observed running \texttt{07\_stability\_test.py} end-to-end: a few minutes of
wall-clock time, with most of the gap between 127 seconds of compute and the longer observed
wall-clock time attributable to one-time model loading and Hugging Face dataset streaming overhead,
not to GPU compute. The calibration-based estimate and the direct observation are consistent, not
off by an order of magnitude, which is the check this section set out to perform.

\section{Standard vs.\ Non-Standard Choices}

\textbf{Standard}: wall-clock timing aggregation and linear extrapolation for compute-budget
estimation is itself standard practice, directly serving efficiency as one of the six metrics
defined for black-box XAI evaluation \citep{arreche2024exai}. Measuring per-sample inference time on
a realistic pilot scale and extrapolating linearly to a deployment-relevant $n$ is the conventional
way to answer "is this affordable at scale," and this chapter follows that convention rather than
departing from it.

\textbf{Non-standard}: the normalized-per-generate-call comparison introduced in Finding 1 is this
chapter's own methodological contribution, not a step drawn from the six-metric framework as
published. Comparing raw per-sample milliseconds across days that bundle different numbers of
underlying \texttt{generate()} calls into a single logged "call" is an easy mistake to make and an
easy one to miss --- Day 6's 215ms genuinely does look cheaper than Day 3's 271ms until the call
count inside each is accounted for, and the naive reading would have flatly contradicted
Chapter~\ref{chap:day06}'s own eager-attention finding. Normalizing to a single-generate-call basis
before drawing any cross-day conclusion is presented here as a detail worth stating explicitly in
any paper's efficiency-methodology paragraph, not a one-off correction specific to this pilot.

\section{Implications for the Paper}

This chapter is the pilot's strong, low-controversy "good news" result. \textbf{Compute is
conclusively not an obstacle to scaling this evaluation framework.} The full six-metric pipeline,
extrapolated linearly from real, mostly-exact per-sample measurements, costs roughly half a GPU-hour
at $n=1{,}000$ on a single consumer-grade RTX 3070 --- well within reach of scaling to the full
3,004-image ConstructionSite test split, and a result that should be reported in the paper without
the qualifications that accompany several of this pilot's other findings. Where Chapters 6, 7, and
others uncovered confounds, instability, or methodological deviations that temper how their numbers
should be read, efficiency's central claim survives those caveats intact: even granting Day 6 and
7's calibration-based estimates the full uncertainty the sanity check in
Section~\ref{sec:day08-gap} affords them, the headline figure would have to be wrong by close to an
order of magnitude before compute became a genuine scaling concern, and the sanity check found no
such discrepancy.

That said, one caveat should be carried forward honestly into Chapters 11 and 12 rather than
omitted for the sake of a clean headline: \textbf{Day 6 and 7's per-sample costs rest on a
15-sample calibration, not the full 163-sample run}, unlike Day 3 and 5's exact aggregation over all
163 samples. This is good enough for the order-of-magnitude story this chapter tells, and the
sanity check in Section~\ref{sec:day08-gap} found it consistent with direct observation, but it
should be re-measured directly --- by instrumenting \texttt{inference\_ms} into Days 6 and 7's
scripts themselves, the same way Day 3 already does --- before quoting precise per-sample numbers
in anything more formal than this pilot. The methodological recommendation worth carrying into the
paper's own efficiency-methodology paragraph is Finding 1's normalization point: any reported
per-sample or per-call timing for a multi-metric XAI pipeline should state explicitly how many
underlying model-generate calls are bundled into that figure, since, as this chapter showed, failing
to do so can make one method look artificially cheaper than another for reasons that have nothing to
do with which method is actually faster.
